\documentclass[11pt]{article}
\usepackage[T1]{fontenc}
\usepackage{lmodern}
\usepackage[margin=1in]{geometry}
\usepackage{amsmath,amssymb,amsthm}
\usepackage{microtype,booktabs,longtable,array}
\usepackage[dvipsnames]{xcolor}
\usepackage[colorlinks=true,linkcolor=MidnightBlue,urlcolor=MidnightBlue,citecolor=MidnightBlue]{hyperref}
\hypersetup{pdftitle={Critical hard-core influence bounds via nonbacktracking collisions},
pdfauthor={Dev D. Goyal},
pdfsubject={Fully autonomous agentic research},
pdfkeywords={hard-core model, spectral independence, criticality, nonbacktracking walk, Lean 4}}
\usepackage{fancyhdr}
\pagestyle{fancy}
\fancyhf{}
\fancyhead[L]{\small Critical hard-core correlations}
\fancyhead[R]{\small Dev D. Goyal}
\fancyfoot[C]{\thepage}
\setlength{\headheight}{14pt}
\setlength{\parskip}{4pt}
\setlength{\parindent}{0pt}
\allowdisplaybreaks
\newtheorem{theorem}{Theorem}[section]
\newtheorem{lemma}[theorem]{Lemma}
\newtheorem{corollary}[theorem]{Corollary}
\newtheorem{proposition}[theorem]{Proposition}
\theoremstyle{definition}
\newtheorem{definition}[theorem]{Definition}
\newcommand{\E}{\mathbb E}
\newcommand{\Prb}{\mathbb P}
\newcommand{\SI}{\operatorname{SI}}
\newcommand{\lamc}{\lambda_{\mathrm c}}
\newcommand{\clog}{\operatorname{clog}}
\newcommand{\abs}[1]{\left|#1\right|}
\newcommand{\code}[1]{\texttt{\detokenize{#1}}}
\newcommand{\replaystatus}{A fresh replay on 30 October 2026, compiling all 32 modules from source against a Mathlib checkout built at the pinned revision, passed; its record is in \code{verification/}.}

\title{\textbf{Critical hard-core influence bounds\\via nonbacktracking collisions}\\[6pt]
\large With a Lean 4 formalization}
\author{Dev D. Goyal}
\date{30 October 2026}

\begin{document}
\maketitle

\section*{Research process and AI disclosure}
This is fully autonomous agentic research. Dev D. Goyal is responsible for publication. No expert validation or external peer review is claimed, and the novelty assessment is provisional.

The research, the proof, the Lean formalization and an earlier version of this note were produced by autonomous AI research agents in an earlier session. The records of that session state that the note was prepared with ChatGPT; the model version is not recorded.

Opus 5.5 by Anthropic used with Claude Code assessed the package for this release. It coordinated an independent AI referee check, which recomputed exact influence row sums by brute force on 57 small graphs at the critical activity, with and without pinnings, checked the collision ingredients on small regular graphs, and compared the Lean statements with the claims. It ran the fresh Lean replay described in Section~\ref{sec:formal} and revised this manuscript: it added the credit to Sinclair, Srivastava, \v{S}tefankovi\v{c} and Yin, the comparison with recent Ising results, the discussion of the constant, and the conditional status of the mixing corollary. AI referee checks are not human peer review. Lean verification establishes only the formal statements described in Section~\ref{sec:formal}.

\begin{abstract}
For each fixed integer $\Delta\ge3$, we establish a uniform
$O_\Delta(\sqrt n\log(n+1))$ bound on the absolute row sums of the hard-core
influence matrix on every finite $n$-vertex simple graph of maximum degree
at most $\Delta$, at all homogeneous fugacities
$0<\lambda\le\lamc(\Delta)=(\Delta-1)^{\Delta-1}/(\Delta-2)^\Delta$.
The bound holds under every feasible pinning. An explicit finite constant
is given below. The proof uses exact vertex-deletion identities, a known
local squared-message contraction, a self-avoiding-walk transfer of the
type introduced by Sinclair, Srivastava, \v{S}tefankovi\v{c} and Yin, and a
finite reconstruction of the recent nonbacktracking collision argument of
Bencs et al. The actual Gibbs
distribution, conditional influences, graph completion, path-count transfer,
and collision estimate are formalized in Lean 4.

The contribution is the resulting general critical hard-core bound;
its ingredients are credited prior mathematics. It recovers, up to a
logarithm, an $O(\sqrt n)$ bound that was stated in 2025 and withdrawn in
January 2026. Through a published proposition that we have not formalized,
it gives a polynomial bound on the mixing time of Glauber dynamics at
criticality. The explicit constant is large.
\end{abstract}

\section{Result and scope}

Let $G=(V,E)$ be a finite simple graph and $n=|V|\ge1$. Its hard-core
distribution at activity $\lambda>0$ is
\[
 \mu_{G,\lambda}(I)=\frac{\lambda^{|I|}}{Z_G(\lambda)},\qquad
 Z_G(\lambda)=\sum_{I\subseteq V\text{ independent}}\lambda^{|I|}.
\]
We use the occupation variables $X_v=\boldsymbol1_{v\in I}$.
A pinning consists of disjoint sets $P,Q\subseteq V$, whose vertices are
required to be occupied and unoccupied respectively. It is feasible exactly
when $P$ is independent. Write $\mu^{P,Q}$ for the resulting conditional law.
For a coordinate $u$ with $0<\mu^{P,Q}(X_u=1)<1$, define
\[
 \Psi^{P,Q}(u,v)=
 \mu^{P,Q}(X_v=1\mid X_u=1)-
 \mu^{P,Q}(X_v=1\mid X_u=0).
\]
This convention has diagonal $1$. Deterministic target coordinates
contribute zero. A uniform absolute-row bound is the
$\ell_\infty$ form of spectral independence; it also controls the usual
matrix eigenvalue notion. The main formal theorem states the absolute-row
inequality directly.

\begin{theorem}[Explicit critical hard-core bound]\label{thm:main}
Fix an integer $\Delta\ge3$, put $d=\Delta-1$, and let $G$ have maximum
degree at most $\Delta$. Suppose $0<\lambda\le d^d/(d-1)^{d+1}$. Set
\[
 N=2\Delta n,\quad r=\lfloor\sqrt N\rfloor+1,\quad
 L=\left\lceil\log_d r\right\rceil,\quad
 M=1+12\Delta d\frac{\lambda}{1+\lambda}\,rL.
\]
For every feasible pinning $(P,Q)$ and every nondegenerate source $u$,
\begin{equation}\label{eq:main}
 \sum_{v\in V}\abs{\Psi^{P,Q}(u,v)}\le M.
\end{equation}
In particular the bound is $O_\Delta(\sqrt n\log(n+1))$, uniformly in
$0<\lambda\le\lamc(\Delta)$ and in the pinning.
\end{theorem}

The degree bound is fixed in the asymptotic notation. The displayed
constant remains explicit when the degree varies. The variable $n$ in
the formal theorem is the original ambient vertex count. No expansion,
girth, bipartiteness, connectivity, or minimum-degree assumption is imposed
on $G$. Activity is homogeneous; arbitrary coordinate-dependent field
tilts are not part of the theorem. The exact $O(\sqrt n)$ hard-core upper
bound and the general critical two-spin question remain unresolved by
this result. The established hard-core lower bound is
$\Omega(n^{1/3})$ for spectral independence \cite{CCYZ}; accordingly,
we make no assertion of optimality up to logarithms.

\section{What is established by the formal development}\label{sec:formal}

All 32 Lean modules compile with the pinned toolchain, and the central
declarations depend only on \code{propext}, \code{Classical.choice} and
\code{Quot.sound}. \replaystatus

The central declarations in \code{CriticalHardcore/Main.lean} are:
\begin{itemize}
\item \code{critical_hardcore_influence_rows};
\item \code{critical_hardcore_spectral_independence};
\item \code{critical_hardcore_pinned_rows}.
\end{itemize}
The last declaration exposes the original pinning
sets, conditional probabilities, and nondegeneracy premise. Its assumptions
are only a finite simple graph, positive activity at or below the stated
threshold, the degree bound, and feasibility/nondegeneracy of the pinning.
There is no assumed tree comparison or collision bound.

The finite Gibbs probabilities are defined as ratios of sums of
$\lambda^{|I|}$ over actual independent sets satisfying the relevant events.
All needed normalizers are proved positive. The formalization does not
define the desired influence in terms of a surrogate recurrence and then
assume that it represents the graph distribution.

\section{Exact finite graph identities}

For a set $S\subseteq V$, use $Z(S)$ and $\E_S$ for the partition function
and expectation of the induced hard-core model on $S$. If $u\in S$, let
$F(S,u)=S\setminus N_G[u]$. Splitting independent sets according to $u$
gives
\begin{equation}\label{eq:split}
 Z(S)=Z(S\setminus\{u\})+\lambda Z(F(S,u)),\qquad
 p(S,u)=\frac{\lambda Z(F(S,u))}{Z(S)}.
\end{equation}
Both terms are positive. Also $0<p(S,u)<1$ and
$p(S,u)\le p_{\max}:=\lambda/(1+\lambda)$: removing additional vertices
can only remove nonnegative summands from the partition function.

For $f:V\to\mathbb R$ define the additive observable
$F_f(I)=\sum_{v\in I}f(v)$ and its conditional difference
\[
 D(S,u;f)=\E_S[F_f\mid X_u=1]-\E_S[F_f\mid X_u=0]
 =f(u)+\E_{F(S,u)}F_f-\E_{S\setminus\{u\}}F_f.
\]
Choose any ordering $v_1,\ldots,v_k$ of the neighbors of $u$ lying in $S$.
Put
\[
 S_i=S\setminus\{u,v_1,\ldots,v_{i-1}\},\qquad
 p_i=p(S_i,v_i).
\]
Each $v_i$ belongs to $S_i$. Successively conditioning $v_i$ to be absent
in these contexts gives the two exact identities
\begin{align}
 p(S,u)&=\frac{\lambda\prod_{i=1}^k(1-p_i)}
                    {1+\lambda\prod_{i=1}^k(1-p_i)},\label{eq:message}\\
 D(S,u;f)&=f(u)-\sum_{i=1}^k p_iD(S_i,v_i;f).\label{eq:observable}
\end{align}
For completeness, $Z(S_i\setminus\{v_i\})/Z(S_i)=1-p_i$ by
\eqref{eq:split}. Multiplying these ratios telescopes from
$S\setminus\{u\}$ to $F(S,u)$ and proves \eqref{eq:message}.
The mixture identity
\[
 \E_{S_i}F_f-\E_{S_i\setminus\{v_i\}}F_f
 =p_i D(S_i,v_i;f)
\]
telescopes over the same contexts and proves \eqref{eq:observable}.
The neighbors need not be independent.

Self-avoiding graph reductions are established methodology \cite{Weitz}.
Iterating these exact identities builds a finite deletion computation tree.
A node stores $(S,u)$ and its actual message $p(S,u)$; its children store
the pairs $(S_i,v_i)$. The free set strictly decreases at each edge, so
fuel $|S|$ exhausts the recursion exactly. Child labels are distinct, and
every descendant free set excludes each of its ancestor labels. Consequently
each tree path labels a self-avoiding graph path, and two distinct tree paths
with the same depth have different full label sequences.

For $|f(v)|\le1$, taking absolute values in \eqref{eq:observable} bounds
$|D(S,u;f)|$ by one plus the sum, over all nonempty downward tree paths,
of the product of their messages, excluding the source root's own message.
This is a finite identity and bound. The computation tree can have far more
than $n$ nodes; no bound on its size by $n$ is used.

Finally,
\begin{equation}\label{eq:duality}
 D(S,u;f)=\sum_{v\in S}f(v)\Psi_S(u,v).
\end{equation}
Taking $f(v)$ to be the sign of $\Psi_S(u,v)$ makes the right side its
absolute row sum. Thus it suffices to bound the weighted tree expansion
uniformly over all bounded observables. This also fixes the row orientation
and the normalization of the influence matrix.

\section{The known local inequality and the path estimate}

The following local inequality is prior art; the stronger squared-message
form is stated in the 2025 Chen--Jiang proceedings and traced there to
Li--Lu--Yin's contraction argument \cite{CJ25,LLY}. The subsequent correction
to the graph comparison does not invalidate this scalar inequality, which
is proved directly here and in Lean.

\begin{lemma}[Local squared-message contraction]\label{lem:local}
Let $d\ge2$, $0\le\lambda\le d^d/(d-1)^{d+1}$, and let
$p_1,\ldots,p_k\in[0,1]$ with $k\le d$. If
$p=\lambda\prod_i(1-p_i)/(1+\lambda\prod_i(1-p_i))$, then
\[
 d p\sum_i p_i\le1.
\]
\end{lemma}
\begin{proof}
Pad the list with zeros to length $d$, without changing its sum or product.
Write $s=\sum_i p_i$ and $A=\prod_i(1-p_i)$. If $ds\le1$, the conclusion
follows from $p\le1$. Otherwise apply AM--GM to the $d+1$ nonnegative
numbers
\[
 1-p_1,\ldots,1-p_d,\quad (ds-1)/d.
\]
Their average is $(d-1)/d$. Therefore
\[
 A(ds-1)\le\frac{(d-1)^{d+1}}{d^d}=\frac1{\lamc(d+1)}.
\]
It follows that $\lambda A(ds-1)\le1$, which rearranges to
$d\lambda As/(1+\lambda A)\le1$.
\end{proof}

Every nonroot node in the deletion tree has at most $d=\Delta-1$ children:
its graph parent has already been deleted. Hence Lemma~\ref{lem:local}
applies at every such node. The original root may have $d+1$ children and
is treated separately.

\begin{lemma}[Weighted paths versus their count]\label{lem:weighted}
Let $d\ge2$, and let a finite rooted tree carry messages $0\le p_x\le p_{\max}$ satisfying
$d p_x\sum_{y\text{ child of }x}p_y\le1$ at every node. Let $a_k$ be
the number of root-to-depth-$k$ paths and let
\[
 M_k=\sum_{(x_0,\ldots,x_k)}\prod_{i=0}^k p_{x_i}.
\]
Then $M_k\le p_{\max}\sqrt{a_k/d^k}$.
\end{lemma}
\begin{proof}
Define
\[
 E_k=\sum_{(x_0,\ldots,x_k)}
       \left(\prod_{i=0}^{k-1}p_{x_i}^2\right)p_{x_k}.
\]
At a subtree rooted at $x$, $E_0(x)=p_x$ and
$E_{k+1}(x)=p_x^2\sum_y E_k(y)$. Induction and the local inequality give
\[
 d^{k+1}E_{k+1}(x)
 \le d p_x^2\sum_y p_y\le p_x.
\]
The sum of squared full path weights is at most $p_{\max}E_k$, because
the last message is at most $p_{\max}$. Cauchy--Schwarz on the finite
list of paths therefore gives
\[
 M_k^2\le a_k p_{\max}E_k
 \le p_{\max}p_{x_0}\frac{a_k}{d^k}
 \le p_{\max}^2\frac{a_k}{d^k}.
\]
All quantities are nonnegative, so taking square roots proves the claim.
No probabilistic independence of tree branches is involved.
\end{proof}

\section{Regular completion and exact walk normalization}

\begin{lemma}[Explicit regular completion]\label{lem:completion}
Every finite graph $G$ of maximum degree at most $\Delta\ge1$ embeds as
an induced subgraph of a simple $\Delta$-regular graph $H$ with exactly
$2\Delta |V(G)|$ vertices.
\end{lemma}
\begin{proof}
Take vertices $(v,a,b)$ with $v\in V(G)$, $a\in\mathbb Z/\Delta\mathbb Z$,
and $b\in\{0,1\}$. Retain a copy of every original edge for each fixed
$(a,b)$. At each $v$, let $\delta_v=\Delta-\deg_G(v)$. For each
$j\in\{0,\ldots,\delta_v-1\}$, join
$(v,a,0)$ to $(v,a+j,1)$ for all $a$.
These are $\delta_v$ disjoint perfect matchings on the copy coordinates.
They have no loops and are disjoint from the original edges, which change
$v$. Thus every new vertex has degree $\deg_G(v)+\delta_v=\Delta$.
Fixing $(a,b)$ gives an induced original copy. The stated vertex count
is immediate.
\end{proof}

Let $H$ now be $\Delta$-regular, $d=\Delta-1$, and fix an oriented edge
$e=(a,b)$. A nonbacktracking walk starts with this exposed edge and then,
at each step, chooses uniformly among the $d$ neighbors different from
its previous vertex. Let
\[
 s_e(k)=\Prb(\text{the initial edge and the next $k$ steps have no repeated vertex}).
\]
With a fixed enumeration of each directed edge's available next neighbors,
there is a bijection between length-$k$ choice words over $\{1,\ldots,d\}$
and these nonbacktracking continuations. Hence
\begin{equation}\label{eq:survival}
 s_e(k)=\frac{\#\{\text{self-avoiding continuations of $e$ with $k$ further steps}\}}{d^k}.
\end{equation}
The sequence is nonincreasing and vanishes once its required number of
distinct vertices exceeds $|V(H)|$.

Apply Lemma~\ref{lem:weighted} to a child subtree of the original source
$u$, whose first child label is $v$. Prepending $u$ to a tree path yields
an injective map into the self-avoiding paths counted by
\eqref{eq:survival}, after embedding $G$ into $H$. Different tree paths
remain different vertex sequences. Thus
\begin{equation}\label{eq:transfer}
 M_k\le p_{\max}\sqrt{s_{(u,v)}(k)},\qquad
 \sum_w|\Psi_G(u,w)|
 \le 1+\sum_{v\text{ a child of }u}p_{\max}\sum_{k\ge0}\sqrt{s_{(u,v)}(k)}.
\end{equation}
The displayed infinite sum is only notation for an eventually zero sequence;
the Lean proof uses uniform bounds on all finite partial sums. The number
of children in the last expression is at most $\Delta$.

\section{A finite square-root survival bound}

This section reconstructs the regular-graph collision method in
Bencs et al.\ \cite{Bencs}. Its purpose is to discharge the complete
walk prerequisite in Lean and obtain the square-root survival sum needed
by \eqref{eq:transfer}. Neither the basic collision phenomenon nor the
endpoint-packing idea is claimed as new here.

\begin{proposition}[Fixed frontier length]\label{prop:collision}
Let $d\ge2$. For every $(d+1)$-regular finite simple graph $H$ with $N$ vertices,
every initial oriented edge $e$, every integer $h\ge0$, and every finite
cutoff $m$,
\begin{equation}\label{eq:collision}
 \sum_{k=0}^{m-1}\sqrt{s_e(k)}
 \le (2h+1)\left(\left\lfloor\frac{2N}{d^h}\right\rfloor+4d^{h+1}\right).
\end{equation}
\end{proposition}

\subsection{Endpoint packing along a surviving history}

Put $\ell=h+1$. At an oriented edge position $i$ of a self-avoiding history,
a frontier is a nonbacktracking path of $\ell$ edges beginning with that
edge. There are $d^h$ such frontiers. Call a frontier safe if, from every
oriented edge along it, no nonbacktracking continuation adding between
one and $\ell$ edges can end at an already visited vertex of the old
history up to position $i+1$.

Safe frontiers at all positions of one self-avoiding history have distinct
endpoints. Here is the full injection argument. Suppose two length-$\ell$
frontiers $a_0,\ldots,a_\ell$ and $b_0,\ldots,b_\ell$ have the same
last vertex, with the second one starting no earlier in the history.
Then $a_0$ is in the second frontier's old history. If
$a_{k+1}=b_{k+1}$ but $a_k\ne b_k$, the walk
\[
 b_k,a_{k+1},a_k,a_{k-1},\ldots,a_0
\]
is nonbacktracking. Its first two vertices are the $k$th oriented edge
of the second frontier, and it adds $k+1\le\ell$ edges to reach the
old history. This contradicts safety. Starting at their equal endpoints
and descending in $k$ shows equality of every vertex of the two frontiers.
Their initial vertices then identify the same history position, because
the history is self-avoiding. The frontiers are identical.

Consequently the total number of safe frontiers over every history position
is at most $N$. Call a position good when at least half of its $d^h$
frontiers are safe. Along any surviving history the number of good positions
is therefore at most
\begin{equation}\label{eq:goodbudget}
 K=\left\lfloor\frac{2N}{d^h}\right\rfloor.
\end{equation}
The same bound applies to any selected trial-start positions.

\subsection{A bad start forces a uniform collision chance}

Fix the past history. A partial choice word is dangerous if some continuation
of length at most $\ell$ makes the complete graph walk repeat a vertex.
If a length-$h$ word has no dangerous prefix, its frontier is safe: a
short return to an old vertex from any frontier edge can be realized as
such a killing continuation from the corresponding prefix. At a bad start,
at least half of the $d^h$ words must therefore encounter danger.

For each dangerous prefix choose one killing continuation and pad it to
exactly $\ell$ choices. Padding is harmless because a collision, once
present, persists under extension. For any dangerous length-$h$ word $w$,
stop at its first dangerous prefix $p$ and encode $w$ as
\[
 p\;\Vert\;\text{chosen padded killing word for }p\;\Vert\;
 \text{the original suffix of }w.
\]
This is a killed word of length $h+\ell=2h+1$. The map is injective:
the first dangerous prefix is still $p$ in the encoded word, since all
earlier prefixes are unchanged; after recovering $p$, delete its fixed
length-$\ell$ killing word to recover the original suffix. More formally,
two minimal dangerous prefixes of a common word are comparable in prefix
order and must coincide.

Thus at least $d^h/2$ words of length $2h+1$ are killed. Conditional on
any bad trial start, the survival probability through the next
$B=2h+1$ steps is at most
\begin{equation}\label{eq:q}
 q=1-\frac{1}{2d^{h+1}}.
\end{equation}
This argument uses counts of complete words, so it does not condition on
a retrospectively selected frontier or assume independence of danger events.

\subsection{Adaptive trials and summation}

Break the continuation into blocks of length $B$. A trial start's good/bad
status can depend on the whole history. Represent all surviving block
extensions as a finite rooted tree with their conditional probabilities.
At a good node the outgoing mass is at most $1$; at a bad node it is
at most $q$. Every live root-to-node path has at most $K$ good starts,
by \eqref{eq:goodbudget}.

Induction on the remaining number of trials and good-start budget proves
\begin{equation}\label{eq:tail}
 s_e(Br)\le q^{(r-K)_+}\qquad(r\ge0).
\end{equation}
At a good node one consumes one unit of budget and at most unit mass;
at a bad node one preserves the budget and multiplies by at most $q$.
An empty extension list contributes zero. These statements apply to every
child separately, so the induction remains valid for history-dependent
statuses. It does not replace adaptive conditional laws by a product of
unconditional probabilities.

By monotonicity of $s_e$ and a geometric sum,
\begin{align*}
 \sum_{k=0}^{m-1}\sqrt{s_e(k)}
 &\le B\sum_{r\ge0}(\sqrt q)^{(r-K)_+}\\
 &=B\left(K+\frac{1}{1-\sqrt q}\right)
 \le B\left(K+\frac{2}{1-q}\right)
 =B\left(K+4d^{h+1}\right).
\end{align*}
In the Lean development the same computation is carried out with finite
sums, avoiding any convergence premise. This proves
Proposition~\ref{prop:collision}.

\subsection{Choosing the square-root scale}

Set $r=\lfloor\sqrt N\rfloor+1$,
$L=\lceil\log_d r\rceil\ge1$, and $h=L-1$. Then
\[
 r\le d^L\le dr,\qquad N<r^2,\qquad
 \left\lfloor\frac{2N}{d^{L-1}}\right\rfloor\le2dr.
\]
Substituting into \eqref{eq:collision} gives
\begin{equation}\label{eq:scale}
 \sum_{k=0}^{m-1}\sqrt{s_e(k)}\le12drL.
\end{equation}
The formal natural-number definitions are
\begin{align*}
 \operatorname{frontierLength}(d,N)&=\clog_d(\lfloor\sqrt N\rfloor+1),\\
 \operatorname{collisionScale}(d,N)&=12d(\lfloor\sqrt N\rfloor+1)
 \operatorname{frontierLength}(d,N).
\end{align*}
Lean also proves the real-variable comparison
\[
 \operatorname{collisionScale}(d,N)
 \le12d(\sqrt N+1)\bigl(1+\log_d(\sqrt N+1)\bigr).
\]
This establishes the stated asymptotic order without suppressing a
graph-dependent constant.

\section{Completing the theorem, including pinnings}

Combine the induced embedding from Lemma~\ref{lem:completion}, the path
transfer \eqref{eq:transfer}, and \eqref{eq:scale}, with $N=2\Delta n$.
There are at most $\Delta$ source children, each contributing at most
$p_{\max}\operatorname{collisionScale}(d,N)$. The source's diagonal term
is $1$. This proves \eqref{eq:main} for every induced free-vertex set
$S\subseteq V$ and every $u\in S$.

Now fix feasible occupied/unoccupied pins $P,Q$. The residual free set is
\[
 R=V\setminus(P\cup Q\cup N_G(P)).
\]
The map $J\mapsto P\cup J$ is a bijection from independent sets of $G[R]$
to independent sets satisfying the original pins. Their weights differ by
the common positive factor $\lambda^{|P|}$, so this map identifies the
residual law exactly with the homogeneous hard-core law on $R$.
The same factor cancels from both additional conditional distributions
used to define an influence. Contributions of the occupied pins to an
additive observable cancel from their difference.

Every vertex in $R$ has a nondegenerate marginal: the empty residual
configuration and its singleton both have positive weight. Vertices outside
$R$ are forced, with marginal either zero or one. Thus the nondegenerate
source coordinates are exactly $R$, and fixed target coordinates contribute
zero. The already proved free-set row inequality on $R$ gives the original
pinned row inequality. All bounds retain the original $n$, so they are
uniform over the pinning. This completes Theorem~\ref{thm:main}.

\section{A conditional consequence for critical Glauber dynamics}

This section derives a mixing bound from a published proposition. The
Markov-chain mixing theorem is not formalized in this package, and we could
not re-examine the current text of that proposition while preparing this
release, so the bound below is conditional on it as quoted.

The corrected Chen--Jiang theorem gives the subcritical spectral-independence
bound $\rho/\delta$, where
\[
 \rho=2+\frac4{\Delta-2},\qquad \lambda=(1-\delta)\lamc(\Delta).
\]
Their Proposition A.1 \cite{CJ26}, based on the homogeneous reweighting
localization framework of \cite{CCYZ}, requires this property for all
pinnings. It does not require arbitrary site-dependent fugacities.
Theorem~\ref{thm:main} supplies an additional uniform cap
$M_n=O_\Delta(\sqrt n\log n)$ for this entire homogeneous family.
Use $K(\delta)=\min\{M_n,\rho/\delta\}$, with $K(0)=M_n$.

The published proposition bounds the mixing time by
\[
 Cn^2\log\Delta\,
 \exp\left(\int_0^\theta\frac{K(\delta)}{1-\delta}\,d\delta\right),
\]
when $(1-\theta)\lamc(\Delta)\le1/(2\Delta)$.
Since $\lamc(\Delta)\le12/\Delta$, the choice $\theta=23/24$ meets this
condition. For large enough $n$, $a=\rho/M_n\le\theta$ and
\begin{align*}
 \int_0^\theta\frac{K(\delta)}{1-\delta}\,d\delta
 &=M_n\log\frac{M_n}{M_n-\rho}
   +\rho\log\frac{\theta(M_n-\rho)}{\rho(1-\theta)}\\
 &\le24\rho+\rho\log\frac{23M_n}{\rho}.
\end{align*}
The finitely many smaller $n$ are absorbed into a degree-dependent constant.
Therefore
\begin{equation}\label{eq:mixing}
 T_{\rm mix}
 =O_\Delta\!\left(
 n^{3+2/(\Delta-2)}(\log n)^{2+4/(\Delta-2)}\right).
\end{equation}
The corrected comparison bound in \cite{CJ26} is
$O(n^{4+4/(\Delta-2)}\log\Delta)$. Equation~\eqref{eq:mixing} improves the
polynomial exponent, conditionally on the quoted proposition. It is not a
claim of optimal critical mixing.

\section{Novelty assessment and significance}

\subsection{The version issue that must not be missed}

The APPROX/RANDOM 2025 proceedings version of Chen--Jiang states a stronger
$O(\sqrt n)$ critical hard-core influence bound \cite{CJ25}. However, the
authors' 7 January 2026 revision explicitly withdraws that claim in
Remark 1.2: the required comparison between general graphs and a truncated
tree quantity was incorrect \cite{CJ26}. The corrected paper retains the
optimal subcritical result and an improved polynomial mixing bound. Search
engines still prominently return the obsolete proceedings statement.

Our proof makes no such truncated-tree comparison. It retains the entire
finite deletion tree, which may be exponentially large, and controls its
weighted paths by counts of self-avoiding graph paths. Every use of the
original graph size comes from actual vertices in a regular completion and
the endpoint-packing argument. This is the specific logical difference from
the withdrawn route.

\subsection{Closest valid sources}

Bencs et al.\ prove a nonbacktracking self-intersection bound and derive an
Ising application \cite{Bencs}. Their Section 5 asks whether critical
two-spin systems on bounded-degree regular graphs always have
$\widetilde O(\sqrt n)$ spectral independence. Theorem~\ref{thm:main}
addresses the hard-core case of that question, and covers maximum-degree
graphs through regular completion. It does not settle the entire
two-spin question.

The step that bounds the weighted tree paths by self-avoiding walk counts
through a potential and the Cauchy--Schwarz inequality follows the
connective-constant method of Sinclair, Srivastava, \v{S}tefankovi\v{c} and
Yin \cite{SSSY}. Pei \cite{Pei} proves an $O(\sqrt n)$ bound without the
logarithm for ferromagnetic Ising models at the uniqueness threshold,
uniformly in fields and pinnings; that result concerns a different model.
Dozier \cite{Dozier} gives an independent birthday-paradox bound for
nonbacktracking walks on regular graphs.

Efthymiou's work on the local connective constant and nonbacktracking
matrices is close methodological prior art \cite{Efthymiou}. Its checked
hard-core statements retain strict subcritical slack. It is important
background for the potential/path-count viewpoint. We make no claim that
the general use of self-avoiding expansions, potentials, or nonbacktracking
counts to bound influence is new.

The 2026 edge-tilting field-dynamics paper \cite{EdgeTilting} gives other
critical mixing results and a spectral-independence estimate with a positive
slack parameter; it does not supply the universal critical finite-size bound
proved here. Peters--Regts--Reppekus \cite{PRR} relate very strong spatial
mixing, zero-freeness, and spectral independence. Their hypothesis is a
uniform decay property of a graph family; it is not an unconditional
critical square-root bound on all bounded-degree graphs.

\subsection{What is and is not proposed as original}

The proposed original mathematical contribution is the universal critical
hard-core $\ell_\infty$ influence bound and its resulting quantitative
consequences. Its proof aligns an exact Gibbs deletion expansion with the
finite square-root survival estimate. The local contraction and the
collision input are prior results, rederived rather than postulated in the
formal development. Formalizing known ingredients is part of verification,
not evidence of their novelty.

The result gives a finite-size critical correlation bound for the
hard-core model on all graphs of maximum degree $\Delta$, under every
feasible pinning, and through existing machinery a better mixing exponent.
Its ingredients are known; the main new observations are that the hard-core
argument needs the sum of square roots of survival probabilities rather than
an expected self-intersection time, and a tail form of the collision bound.
The constant is large: the bound improves on the trivial bound $n$ only for
$n\ge3.4\times10^{6}$ when $\Delta=3$ and $n\ge4.2\times10^{6}$ when
$\Delta=4$. Monte Carlo estimates on random cubic graphs with up to $10^5$
vertices suggest that the $\sqrt n$ order of the survival sum is real and that
the proved constant is roughly 40--70 times too large; these are numerical
observations, not theorems.

The post-result review checked primary papers, corrected versions,
conference proceedings, and related publications, with broader discovery
searches over preprints, institutional repositories, thesis records,
bibliographic indexes, MathOverflow, and GitHub. Search terms included
hard-core/hardcore, critical/uniqueness threshold, total influence,
$\ell_\infty$ spectral independence, self-avoiding paths, nonbacktracking
walks, and square-root finite-size bounds. No valid earlier theorem with
the same general conclusion was located in the reviewed sources. Searches
cannot exclude unpublished work or guarantee exhaustive indexing. The
accessible bibliographic coverage included DBLP; direct access to zbMATH
and MathSciNet did not succeed, and some thesis PDFs were unavailable.
The GitHub search is likewise not exhaustive. No priority
claim is made beyond this documented search.

\section{Formal correspondence and reproducibility}

\subsection{Source map}

\begin{longtable}{>{\raggedright\arraybackslash}p{0.38\linewidth}>{\raggedright\arraybackslash}p{0.54\linewidth}}
\toprule
Files under \code{lean/CriticalHardcore/} & Mathematical role \\
\midrule
\endhead
\code{GibbsModel.lean}, \code{GibbsInfluence.lean} & Actual weighted independent-set sums, positive normalizers, conditional probabilities, and row-sum duality.\\
\code{DeletionRecursion.lean}, \code{ComputationTree.lean}, \code{ComputationGeometry.lean} & Exact neighbor deletion, terminating computation tree, valid states, and self-avoiding labels.\\
\code{LocalContraction.lean}, \code{WeightedTree.lean}, \code{WeightedTotals.lean}, \code{WeightedPaths.lean} & AM--GM contraction, second-moment energy, weighted path sums, and sequence injectivity.\\
\code{PathCountBridge.lean}, \code{InfluenceTransfer.lean} & Injection into actual graph path counts and transfer back to actual Gibbs influences.\\
\code{RegularCompletion.lean} & Explicit induced embedding into a $\Delta$-regular graph on $2\Delta n$ vertices.\\
\code{Nonbacktracking*.lean} & Actual graph walks, uniform choice-word bijections, endpoint packing, local collision, normalized survival, and adaptive trials.\\
\code{FirstDanger.lean}, \code{AdaptiveTrials.lean}, \code{WordTrial*.lean} & Injective killing-word encoding and finite history-dependent budget induction.\\
\code{CollisionScale.lean}, \code{CollisionBound.lean}, \code{ScaleInterpretation.lean} & Explicit square-root/logarithm scale and the completed actual-graph walk bound.\\
\code{Pinning.lean}, \code{SpectralStatement.lean}, \code{Main.lean}, \code{Audit.lean} & Actual pinning laws, nondegenerate coordinates, central theorem, and executable statement/axiom audit.\\
\bottomrule
\end{longtable}

\subsection{Pinned environment and clean build}

The compiler is Lean 4.19.0, commit \code{6caaee842e94}. Mathlib is pinned to
\begin{center}\small\code{c44e0c8ee63ca166450922a373c7409c5d26b00b}.\end{center}
The included \code{lake-manifest.json} also pins all eight transitive
packages. The verification script checks the compiler version, all package
commits, and clean tracked dependency sources. It copies the project proofs
into a fresh directory and rebuilds every local module in dependency order
with \code{autoImplicit=false}. Its import path excludes all previously
built project objects. Pinned compiled dependency caches are reused
read-only; this is not a claim to rebuild the Lean compiler and all of
Mathlib from source in the present session.

From a clean checkout of the delivered \code{lean/} directory on Linux or
WSL, with Elan, Python 3, Git, curl, tar, and network access available, run:
\begin{verbatim}
lake exe cache get
python3 verify.py
\end{verbatim}
The toolchain file selects Lean 4.19.0, and the manifest records exact
dependency revisions. For an already downloaded pinned cache, run instead:
\begin{verbatim}
python3 verify.py --lean /path/to/lean-4.19.0/bin/lean \
  --mathlib /path/to/pinned/mathlib4
\end{verbatim}
The verifier creates three audit files next to the project:
\begin{itemize}
\item \code{verification/build.log};
\item \code{verification/axioms.txt};
\item \code{verification/certificate.json}.
\end{itemize}
The certificate records every module result, source
SHA-256 digests, dependency commits, and central axiom dependencies.
The included certificate is evidence for the delivered source; rerunning
the script produces a new certificate for the local copy.

No \code{sorry}, \code{admit}, custom \code{axiom}, unchecked native
decision shortcut, or trust-level override is allowed. The only permitted
axioms are the standard Lean foundations \code{propext},
\code{Classical.choice}, and \code{Quot.sound}. Their names are reported
explicitly rather than hidden behind an assertion of being ``axiom-free.''

\subsection{Independent checks and remaining uncertainty}

The referee check described in the disclosure covered root versus
nonroot degree, path multiplicity, recursion termination, positive
conditioning events, forced coordinates, oriented-edge lengths, adaptive
trial dependence, and the regular completion's exact size, and found no
error. These are AI checks, not external peer review. The central inequality
and its square-root/logarithm interpretation are the formal claims. The
mixing consequence imports published analysis and is not independently
Lean formalized. The originality assessment remains bounded by the
literature coverage documented here. The proof does not establish optimal
hard-core scaling, arbitrary inhomogeneous activity, or the full critical
two-spin conjecture.

\begingroup
\small
\setlength{\parskip}{0pt}
\begin{thebibliography}{99}
\setlength{\itemsep}{4pt}
\bibitem{Bencs}
F. Bencs, L. A. Goldberg, M. Jenssen, M. Jerrum, G. Pete, G. Regts,
and Y. Yin.
\emph{On the self-intersection time of non-backtracking random walks}.
arXiv:2608.09729v1, 10 August 2026.
\url{https://arxiv.org/abs/2608.09729}.

\bibitem{CJ26}
Z. Chen and T. Jiang.
\emph{Improved Mixing of Critical Hardcore Model}.
arXiv:2505.07515v2, 7 January 2026. In particular Remark 1.2,
the corrected subcritical theorem, and Proposition A.1.
\url{https://arxiv.org/html/2505.07515v2}.

\bibitem{CJ25}
Z. Chen and T. Jiang.
\emph{Improved Mixing of Critical Hardcore Model}.
APPROX/RANDOM 2025, LIPIcs 353, Article 51.
\url{https://doi.org/10.4230/LIPIcs.APPROX/RANDOM.2025.51}.
The critical $O(\sqrt n)$ conclusion in this version was withdrawn in
\cite{CJ26}; the local contraction is independently checked in this note.

\bibitem{CCYZ}
X. Chen, Z. Chen, Y. Yin, and X. Zhang.
\emph{Rapid Mixing at the Uniqueness Threshold}.
arXiv:2411.03413v3, 7 January 2026; STOC 2025 proceedings version.
Use the corrected arXiv version for current claims.
\url{https://arxiv.org/abs/2411.03413}.

\bibitem{LLY}
L. Li, P. Lu, and Y. Yin.
\emph{Correlation Decay up to Uniqueness in Spin Systems}.
SODA 2013, pp. 67--84; revised full version arXiv:1111.7064v3.
\url{https://arxiv.org/abs/1111.7064}.

\bibitem{Efthymiou}
C. Efthymiou.
\emph{On sampling two spin models using the local connective constant}.
arXiv:2411.08179v2; SODA 2026.
\url{https://arxiv.org/html/2411.08179v2}.
\url{https://doi.org/10.1137/1.9781611978971.39}.

\bibitem{EdgeTilting}
X. Chen, Z. Ju, T. Miao, Y. Yin, and X. Zhang.
\emph{Edge-Tilting Field Dynamics: Rapid Mixing at the Uniqueness
Threshold and Optimal Mixing for Swendsen-Wang Dynamics}.
arXiv:2604.10525v2, May 2026.
\url{https://arxiv.org/abs/2604.10525}.

\bibitem{PRR}
H. Peters, G. Regts, and J. Reppekus.
\emph{Decay of correlations and zeros for the hard-core model}.
arXiv:2603.17858v2, 15 June 2026.
\url{https://arxiv.org/html/2603.17858v2}.

\bibitem{SSSY}
A. Sinclair, P. Srivastava, D. \v{S}tefankovi\v{c}, and Y. Yin.
\emph{Spatial mixing and the connective constant: optimal bounds}.
Probability Theory and Related Fields 168 (2017), 153--197.

\bibitem{Pei}
Y. R. Pei.
\emph{Critical and near-critical influence bounds for ferromagnetic Ising models}.
arXiv:2609.20866, September 2026.
\url{https://arxiv.org/abs/2609.20866}.

\bibitem{Dozier}
B. Dozier.
\emph{The Birthday Paradox for non-backtracking walks on regular graphs}.
arXiv:2608.26321, August 2026.
\url{https://arxiv.org/abs/2608.26321}.

\bibitem{Weitz}
D. Weitz.
\emph{Counting independent sets up to the tree threshold}.
STOC 2006, pp. 140--149.
\url{https://doi.org/10.1145/1132516.1132538}.

\bibitem{ALO}
N. Anari, K. Liu, and S. Oveis Gharan.
\emph{Spectral Independence in High-Dimensional Expanders and
Applications to the Hardcore Model}.
FOCS 2020; SIAM Journal on Computing.
\url{https://doi.org/10.1137/20M1367696}.
\end{thebibliography}
\endgroup
\end{document}
